\section{The local factor and conductor exponent at 2}\label{sec:local-two}
Let $\ell\ne2$ and retain the rational descents defined by the literal
cubic deck operation and inversion. On
$T=\PP^1_z\setminus\{0,\infty,\mu_3\}$ let
$f:\mathcal E\to T$ be the family $y^2=x(x-1)(x-z^3)$, and put
\[
 \mathcal H_\ell=H^1_!(T_{\overline\Q},\Sym^4R^1f_*\Q_\ell),
 \qquad W_{4,\ell}^{\pm}=(e_{\nt}\mathcal H_\ell)^{I=\pm1}.
\]
The following argument uses one arithmetic Galois cover for the complete
calculation. The semistable-model criterion supplies tameness; the main,
tail and node terms then determine the inertia invariants.

\begin{theorem}[Dyadic local data]\label{p2:thm:dyadic-local-data}
For every coefficient prime $\ell\ne2$,
\[
 \mathcal H_\ell^{I_2}=0,\qquad
 \Sw_2(\mathcal H_\ell)=0.
\]
Consequently both specified rank-four descents satisfy
\[
 d_{2,\ell}^{\pm}=0,\qquad P_{2,\ell}^{\pm}(T)=1,\qquad
 \Sw_2(W_{4,\ell}^{\pm})=0,\qquad a_2(W_{4,\ell}^{\pm})=4.
\]
\end{theorem}
The proof occupies the remainder of this section. It retains the
arithmetic action on the cover and on its coefficient system separately.
The identity $W^-_{4,\ell}=W^+_{4,\ell}\otimes\chi_{-3}$ is consistent
with the equal dyadic data: $\chi_{-3}$ is unramified at $2$ and has
Frobenius value $-1$ there.

\subsection{The exact Galois closure and its arithmetic action}
Fix an auxiliary prime $\ell\ne2$, put $K_0=\Q_2(\zeta)$, where
$\zeta^3=1\ne\zeta$, and let $L=K_0(\pi)$ with $\pi^3=2$.
Choose an arithmetic inertia lift $\sigma$ of the generator of
$\operatorname{Gal}(L/K_0)$ with $\sigma(\pi)=\zeta\pi$.
The following computations concern actual geometric deck transformations
and this arithmetic automorphism separately.

\subsubsection{The constant in the cubic closure}
Set $h=v^3$, $D=2h+1$, and $z=v(h+2)/D$.  The polynomial
\[
 f_z(X)=X^4-2zX^3+2X-z
\]
factors after adjoining $v$ as
\[
 f_z(X)=(X-v)\left(X^3-\frac{3v}{D}X^2
                 -\frac{3v^2}{D}X+\frac{h+2}{D}\right).
\]
Adjoin $\eta$ satisfying
\begin{equation}\label{p2:eq:closure-constant}
 \eta^3=4(v^3+1).
\end{equation}
The other three roots are exactly
\begin{equation}\label{p2:eq:closure-roots}
 v_i=\frac{v-\zeta^i\eta^2/2-\zeta^{-i}v^2\eta}{D},
 \qquad i=0,1,2.
\end{equation}
Indeed, after the substitution $X=Y+v/D$, the cubic has coefficients
\[
 p=-\frac{6v^2(h+1)}{D^2},\qquad
 q=\frac{2(h+1)(2h^2+h+1)}{D^3}.
\]
The two Cardano terms $u=-\eta^2/(2D)$ and $w=-v^2\eta/D$
satisfy $uw=-p/3$ and $u^3+w^3=-q$.
The discriminant is
\[
 \operatorname{disc}(f_z)=-432(z^3-1)^2
                         =-3\bigl(12(z^3-1)\bigr)^2.
\]
Consequently $K_0(v,\eta)$ is the splitting field over $K_0(z)$:
it contains the four roots and has degree $4\cdot3=12$, since
$v^3+1$ has simple zeros and is not a cube in $K_0(v)$.
Its group is $A_4$.  Over $\Q_2(z)$ the group is $S_4$ and the
constant extension is exactly $K_0$.  For example, the smooth cubic
has the rational point $(v,\eta)=(-1,0)$, so its function field is
geometrically integral and introduces no additional constants.

The smooth projective closure is an elliptic curve, with the displayed
point as origin.  Explicitly,
\[
 x=\frac{3\eta}{v+1},\qquad y=\frac{9(1-v)}{v+1}
 \quad\Longrightarrow\quad y^2=x^3-27.
\]
The normal Klein four subgroup of $A_4$ consists of translations by
the four two-torsion points of this elliptic curve.  The quotient
generator is
\[
 c:(v,\eta)\longmapsto(v,\zeta\eta).
\]
One can verify the translation assertion directly: translation by
$(3,0)$ sends $x$ to $3(x+6)/(x-3)$ and $y$ to
$-27y/(x-3)^2$; the resulting $v$ is $v_0$ in
\eqref{p2:eq:closure-roots}.  Conjugation by $c$ gives the other two
nontrivial translations.

\subsubsection{Why the deck action is the projective three-torsion action}
For the Legendre family with $\lambda=z^3$, put
\[
 \alpha(v)=\frac{(v^3+2)^2}{3(2v^3+1)}.
\]
Direct substitution gives
\[
 3\alpha^4-4(1+\lambda)\alpha^3+6\lambda\alpha^2-\lambda^2=0.
\]
Thus $\alpha$ is the $x$-coordinate of a nonzero point of order three,
or equivalently one of the four lines in the two-dimensional
$\mathbf F_3$-space of three-torsion.  The rational map $v\mapsto\alpha$
has degree six, whereas $v\mapsto z$ has degree four.  Therefore
$\alpha\notin\overline\Q_2(z)$: otherwise the degree six would be a
multiple of four.  The subgroup $C_3$ fixing $v$ is maximal in $A_4$,
so this already implies
\[
 \overline\Q_2(z,\alpha)=\overline\Q_2(v).
\]
The normal closures are equal.  Hence the $A_4$ cover just constructed
is precisely the splitting cover of the four projective three-torsion
lines, with their faithful permutation action.  Quadratic twists do
not change these lines and act trivially on $\operatorname{Sym}^4$.

\subsection{The semistable marked model of the true splitting cover}
Use the splitting cover and notation of the preceding subsection:
\[
\eta^3=4(v^3+1),\qquad e=\eta/\pi^2,\qquad
e^3=v^3+1,\qquad \pi^3=2,
\]
over $L=\Q_2(\zeta_3,\pi)$. Write $h=v^3$.
For completeness, the coordinate changes are explicit. Put
\[
 q=\frac{v^3+2}{v^3-1},\qquad
 \lambda=z^3=\frac{q^3(q+2)}{(q-1)(q+1)^3},\qquad
 \alpha=\frac{q^2}{q^2-1}
       =\frac{(v^3+2)^2}{3(2v^3+1)}.
\]
After quadratically twisting the Legendre curve by $\alpha$, the rational
three-torsion point and its tangent slope are
\[
 X_0=\alpha^2,\qquad
 Y_0=\frac{q^5}{(q-1)^3(q+1)^4},\qquad
 M_0=\frac{q(q^2+q+1)}{(q-1)(q+1)^2}.
\]
Translation by this flex gives $y^2+2M_0xy+2Y_0y=x^3$.
Scale first by $q/((q-1)(q+1))$ and then by $\pi$.
The resulting equation is
\begin{equation}\label{p2:eq:complete-tate}
y^2+\pi^2 A(v)xy+B(v)y=x^3,
\qquad
A=\frac{3(h^2+h+1)}{(h-1)(2h+1)},\quad
B=\frac{(h+2)^2}{(h-1)(2h+1)}.
\end{equation}
The quadratic twist disappears on $\Sym^4R^1$. Throughout this section,
scaling by $k$ means $x=k^2x'$, $y=k^3y'$, so that $(a_1,a_3)$ changes
to $(a_1/k,a_3/k^3)$.

The discriminant before pulling back to $e$ is
\[
\Delta=\frac{3^6h^2(h+1)^2(h+2)^6}{(h-1)^6(2h+1)^6}.
\]
On the $v$-line, the fourteen cusp fibres are $I_6$ at $v=0,\infty$ and
at the three roots of each of $h=-2$, $2h+1=0$, and $h=1$; the three
fibres at $h=-1$ are $I_2$. At either pole divisor $h=1$ or $2h+1=0$,
scaling by the inverse of its local parameter makes $a_1$ a unit,
$a_3$ vanish to order two, and the discriminant have order six.
Thus these poles do not conceal additive fibres.
The cubic cover $e^3=h+1$ ramifies at the three $I_2$ points, changing
them to $I_6$, and is unramified at all other cusps. Consequently all
36 cusp fibres on the true splitting cover are $I_6$.

The projective Fermat curve $C:e^3=v^3+1$ has good reduction at 2.
Its marked cusp sections form exactly nine clusters in the special fibre:
three over $v=0$, three at infinity, and three over $e=0$.
Use respectively $v$, $w=1/v$, and $e$ as local parameters. In each
cluster the four sections have parameter zero and three parameters whose
cubes are $\pm2$. The complementary coordinate is a unique power series
by Hensel's lemma, its derivative being three times a unit.
Blow up the nine disjoint ideals $(\text{parameter},\pi)$.
The resulting proper regular model $\mathcal B/\mathcal O_L$ has one
smooth genus-one main component and nine rational tails, all of
multiplicity one. Every tail has four disjoint cusp sections and meets
the main component at a fifth point. All 36 sections meet the smooth
locus of the special fibre and avoid the nine nodes.

This is the stable model of the curve with its full cusp divisor: each
rational tail has five special points. Uniqueness of the stable marked
model extends the deck group $G=A_4$ and the semilinear arithmetic action
$\sigma$ to $\mathcal B$. It also shows that the explicit blowups do not
depend on choices of labels within a cusp orbit.

\subsubsection{Good models over the nine nodes}
The following equations are obtained directly from
\eqref{p2:eq:complete-tate}. Their denominators are units in the indicated
node neighbourhoods.

At each node over zero put $vr=\pi$ and scale by $k=v^2$. Then
\[
a_1=r^2A(v),\qquad
a_3=\frac{(1+r^3)^2}{(v^3-1)(2v^3+1)}.
\]
At each infinity node put $w=1/v$, $wr=\pi$, and scale by $k=1/w$.
The coefficients become
\[
a_1=\frac{3r^2(1+w^3+w^6)}{(1-w^3)(1+r^3)},\qquad
a_3=\frac{(1+2w^3)^2}{(1-w^3)(1+r^3)}.
\]
At each node over $e=0$ put $er=\pi$, $h=e^3-1$, and scale by $k=1/e$.
Since $h-1=e^3-2=e^3(1-r^3)$, the coefficients are
\[
a_1=\frac{3r^2(h^2+h+1)}{(1-r^3)(2h+1)},\qquad
a_3=\frac{(h+2)^2}{(1-r^3)(2h+1)}.
\]
At the closed point of every node they reduce to $(a_1,a_3)=(0,1)$.
The Tate discriminant $a_3^3(a_1^3-27a_3)$ is therefore a unit.
These are smooth elliptic models over the node rings themselves, not
only good models over the generic points of the components.

\subsubsection{Integral lifts on all three kinds of tails}
Write $U=u^3$. At a zero tail put $v=\pi u$ and scale
\eqref{p2:eq:complete-tate} by $k=\pi^2$. One obtains
\[
a_1=\frac{3(4U^2+2U+1)}{(2U-1)(4U+1)},\qquad
a_3=\frac{(U+1)^2}{(2U-1)(4U+1)}.
\]
At an infinity tail put $w=\pi u$ and scale by
$k=1/(\pi(U+1))$. The integral lift is
\[
a_1=\frac{3(1+2U+4U^2)}{1-2U},\qquad
a_3=\frac{(1+4U)^2(U+1)^2}{1-2U}.
\]
At an $e=0$ tail put $e=\pi u$, so $h=2U-1$, and scale by
$k=1/(\pi(U-1))$. This gives
\[
a_1=\frac{3(h^2+h+1)}{2h+1},\qquad
a_3=\frac{(U-1)^2(h+2)^2}{2h+1}.
\]
Every displayed denominator is a unit along its finite special tail.
All three lifts reduce to the same ordinary elliptic family
\begin{equation}\label{p2:eq:common-tail}
E_u:\quad y^2+xy+(u^3+1)^2y=x^3,
\qquad \Delta=u^6(u^3+1)^6.
\end{equation}
It has four $I_6$ fibres, at $u=0$ and the three points $u^3=1$;
their $a_1$ is a unit. At the attaching point $u=\infty$, scaling by
$u^2$ gives good reduction, agreeing with the node charts.

On the main component, away from its nodes, scale
\eqref{p2:eq:complete-tate} by $v^2/e$. Its reduction is the constant
supersingular curve $E_0:y^2+y=x^3$. The node charts extend this good
reduction at all attaching points. The nine tails account for exactly
$9\cdot4=36$ fibres of type $I_6$, agreeing with the full generic cusp
divisor. Thus no horizontal collision or extra discriminant component
remains. The ratios of the main scale to the three node scales are
$1/e$, $1/(we)$, and $v^2$. The ratios of the three tail scales to
their node scales are $u^{-2}$, $u/(u^3+1)$, and $u/(u^3-1)$.
All are units on their respective overlaps; the latter two node charts
exclude the vanishing of their displayed denominators. These explicit
coordinate transitions glue the Weierstrass line bundle and cubic models.
At a horizontal cusp, $c_4$ is a unit and the exact discriminant divisor
is $\Delta=\text{unit}\cdot t^6$, with $t$ a local equation of its
section. The nodal deformation is therefore etale locally
$xy=\text{unit}\cdot t^6$; the unit can be absorbed in one coordinate.
The resulting
coefficient system has good reduction at every vertical node and
multiplicative reduction precisely along the 36 disjoint horizontal
sections. The functorial proper semistable/logarithmic cohomology
construction is consequently equivariant for $G$ and $\sigma$.

\subsection{From the charts to arithmetic tameness}

\begin{lemma}[Tameness from a semistable marked elliptic model]
\label{p2:lem:p2-log-bridge}
Let $K$ be a finite extension of $\mathbf Q_p$, let $L/K$ be a finite
 tamely ramified extension, and fix a prime $\ell\ne p$.  Let $U/K$ be
 a smooth geometrically connected curve, let $f:E\to U$ be an elliptic
 scheme, and put $\mathcal V=R^1f_*\mathbf Q_\ell$.
Suppose that there is a finite \'etale morphism $g:U'\to U_L$ of degree
 $d>0$ and an elliptic scheme $f':E'\to U'$ with
\[
 \operatorname{Sym}^4R^1f'_*\mathbf Q_\ell
   \simeq g^*\operatorname{Sym}^4\mathcal V.
\]
Suppose in addition that $U'$ has a proper semistable marked model
 $(\mathcal C,\mathcal D)/\mathcal O_L$ with the following properties.
\begin{enumerate}
\item The reduced horizontal divisor $\mathcal D$ is a disjoint union
 of \'etale sections in the smooth locus of $\mathcal C$, and
 $U'=\mathcal C_L\setminus\mathcal D_L$.
\item The elliptic scheme $E'$ extends to a proper flat family
 $\mathcal E\to\mathcal C$ which is smooth over
 $\mathcal C\setminus\mathcal D$.  In particular it is smooth at
 every node of the vertical fibre of $\mathcal C$.
\item Along $\mathcal D$ the family is semistable multiplicative:
 \'etale locally its singularities have equations
 $xy=t^m$ for some $m\ge1$, where $t=0$ defines a section of
 $\mathcal D$.  An invertible factor on the right is allowed.
\end{enumerate}
Then wild inertia acts trivially on
 $H^1(U_{\overline K},\operatorname{Sym}^4\mathcal V)$, and hence
 on its parabolic subspace
\[
 H^1_!(U_{\overline K},\operatorname{Sym}^4\mathcal V)
 :=\operatorname{im}\bigl(H^1_c(U_{\overline K},
 \operatorname{Sym}^4\mathcal V)\longrightarrow
 H^1(U_{\overline K},\operatorname{Sym}^4\mathcal V)\bigr).
\]
The same conclusion holds for every Galois-stable direct summand or
 subquotient of these representations.
\end{lemma}

\begin{proof}
Give $S=\operatorname{Spec}\mathcal O_L$ its standard divisorial log
 structure and $\mathcal C$ the log structure defined by its reduced
 vertical fibre together with $\mathcal D$.  This is a fine and
 saturated log-smooth curve over $S$.  Give $\mathcal E$ the log
 structure defined by the inverse image of this boundary, taking
 its fine and saturated toric model where necessary.  The stated
 node charts show that $\mathcal E\to\mathcal C$ is log smooth.
Here no condition $p\nmid m$ is needed for $xy=t^m$: the toric
 relation is $e_x+e_y=m e_t$, and the group cokernel relative to
 $\mathbf Z e_t$ is free, generated by $e_x=-e_y$.
An invertible factor can be absorbed in $x$ or $y$.  At a vertical
 base node the elliptic family is smooth, so its log chart is just
 the smooth pullback of $uv=\varpi$.  Thus composition gives a
 proper fine and saturated log-smooth model over $S$.

Form the fourfold fibre product in the category of fine and
 saturated log schemes,
\[
 \mathcal X=\mathcal E\mathbin{\times^{\mathrm{fs}}_{\mathcal C}}
 \mathcal E\mathbin{\times^{\mathrm{fs}}_{\mathcal C}}
 \mathcal E\mathbin{\times^{\mathrm{fs}}_{\mathcal C}}\mathcal E.
\]
Log smoothness is stable under this base change and composition.
The underlying scheme of $\mathcal X$ is proper over $S$: its
 toric saturation is finite over the ordinary proper fibre product.
The locus where the log structure is trivial is exactly
\[
 X^\circ=E'\times_{U'}E'\times_{U'}E'\times_{U'}E'.
\]
In particular it is an open subset of the generic fibre. The
 horizontal-boundary theorem \cite[Corollary~8.4.3, p.~310]{Illusie}
 applies to the trivial-log open of this proper log-smooth fine and
 saturated model. The original nearby-cycle and comparison statements are
 \cite[Variant~3.4.1, Proposition~4.2 and Remark~4.2.1]{Nakayama}.
Consequently wild inertia of $L$ is trivial on
 $H^q(X^\circ_{\overline L},\mathbf Z/\ell^n)$ for every $q,n$,
 and hence on its $\mathbf Q_\ell$ cohomology.  One must use the
 horizontal-boundary statement here; the generic log structure on
 the proper compactification is not trivial.

For completeness, the required coefficient cohomology is an
 explicit summand of ordinary degree-five cohomology of
 $X^\circ$.  Write $F:X^\circ\to U'$ and
 $\mathcal V'=R^1f'_*\mathbf Q_\ell$.  Let $\iota_i$ be inversion
 on factor $i$, and let $S_4$ permute the four factors.  The commuting
 idempotents
\[
 e_{\rm odd}=\prod_{i=1}^4\frac{1-\iota_i^*}{2},\qquad
 e_{\rm sym}=\frac1{4!}\sum_{\sigma\in S_4}
                 \operatorname{sgn}(\sigma)\sigma^*,\qquad
 e=e_{\rm odd}e_{\rm sym}
\]
are defined over $L$.  In relative K\"unneth cohomology,
 $e_{\rm odd}$ retains only the tensor of four degree-one factors.
Thus it kills $R^bF_*\mathbf Q_\ell$ for $b\ne4$ and has image
 $(\mathcal V')^{\otimes4}$ for $b=4$.  A transposition of two
 degree-one factors contributes the Koszul sign $-1$; the sign in
 $e_{\rm sym}$ therefore makes its action on this tensor power
 the ordinary symmetrizer.  It follows that
\[
 eR^bF_*\mathbf Q_\ell=
 \begin{cases}
 \operatorname{Sym}^4\mathcal V',&b=4,\\
 0,&b\ne4.
 \end{cases}
\]
Applying this projector to the Leray spectral sequence leaves
 only this one row.  Its canonical edge identification is
\[
 eH^5(X^\circ_{\overline L},\mathbf Q_\ell)
   \simeq H^1(U'_{\overline L},
                    \operatorname{Sym}^4\mathcal V').
\]
All maps are Galois equivariant, so the right-hand side is tame.

Finally the finite \'etale pullback and trace maps satisfy
\[
 \operatorname{Tr}_g\circ g^*=d\,\mathrm{id}
 \quad\text{on}\quad
 H^1(U_{\overline K},\operatorname{Sym}^4\mathcal V).
\]
After dividing the trace by $d$, pullback is a split injection
 into the tame coefficient cohomology just obtained.  The
 parabolic image is a Galois-stable subspace of ordinary
 cohomology.  Since $L/K$ is tamely ramified, its wild inertia
 subgroup equals the wild inertia subgroup of $K$ under the
 usual embedding $G_L\subset G_K$.  This proves all assertions.
\end{proof}


\begin{corollary}[Wild inertia on the cubic Legendre group]
\label{p2:cor:true-cover-swan}
For every $\ell\ne2$, the full parabolic group $\mathcal H_\ell$ and
both specified descents $W_{4,\ell}^{\pm}$ have zero Swan conductor at $2$.
\end{corollary}
\begin{proof}
Apply Lemma~\ref{p2:lem:p2-log-bridge} to the true splitting cover
\[
 z=\frac{v(v^3+2)}{2v^3+1},\qquad e^3=v^3+1,
 \qquad L=\Q_2(\zeta_3,\pi),\quad \pi^3=2.
\]
This constant extension has ramification index three and is tame.
If $h=v^3$, then
\[
 z^3-1=\frac{(h-1)^3(h+1)}{(2h+1)^3},\qquad
 \frac{dz}{dv}=\frac{2(v^3-1)^2}{(2v^3+1)^2}.
\]
The degree-four rational map ramifies only over the original cusps.
The Fermat projection ramifies at $v^3=-1$, which also maps into
$z^3=1$; it is unramified at the three points above infinity.
Removing all preimages of $\{0,\infty,\mu_3\}$ therefore gives a finite
etale degree-twelve cover. Its compactification has genus one and
its removed divisor has the 36 points already accounted for above.
The displayed integral charts and gluing maps verify each model
hypothesis of the lemma: separated etale horizontal sections, good
elliptic reduction at the nine vertical nodes, and horizontal
singularities $xy=\text{unit}\cdot t^6$. The twist parameter $\alpha=(h+2)^2/(3(2h+1))$ is invertible away
from these cusps, so it introduces no additional singular locus on
the generic open. The fourth symmetric power removes this quadratic
twist. The lemma proves tameness
of the full ordinary coefficient cohomology, hence of its parabolic
subspace and the two rational projected summands.
\end{proof}
\subsection{Geometric deck action and arithmetic inertia}
\subsubsection{The good main component and two different order-three actions}
Set
\[
 e=\eta/\pi^2,\qquad e^3=v^3+1.
\]
This Fermat model is smooth and proper over $\mathcal O_L$ and has
supersingular genus-one special fibre $\overline C$.  In particular
$\overline C[2](\overline{\mathbf F}_2)$ has just its origin.
The translations by two-torsion extend to this good model and
specialize to the identity.  Thus $K_4$ fixes the entire main
component pointwise.  In contrast, both $c$ and $\sigma$ act on its
base by $e\mapsto\zeta e$; for $\sigma$ this follows from
$\sigma(\eta)=\eta$ and $\sigma(\pi^2)=\zeta^2\pi^2$.

Their actions on the elliptic coefficient are different.  The exact
Tate equation derived above is
\[
 E^*: y_*^2+\pi^2 A(v)x_*y_*+B(v)y_*=x_*^3,
\quad
 A=\frac{3(h^2+h+1)}{(h-1)(2h+1)},\quad
 B=\frac{(h+2)^2}{(h-1)(2h+1)}.
\]
It was obtained by a scale $\pi$ after transformations defined over
$K_0(v)$, up to the harmless quadratic twist.  Hence
\[
 c:(x_*,y_*)\mapsto(x_*,y_*),\qquad
 \sigma:(x_*,y_*)\mapsto(\zeta x_*,y_*).
\]
On the main component use the further scale $k=v^2/e$, and write
$X=x_*/k^2$, $Y=y_*/k^3$.  Its special equation is the constant curve
\[
 E_0:Y^2+Y=X^3.
\]
For instance its $a_3$ before reduction is
$B(h+1)/h^2$, which reduces to one.  Transport of coordinates gives
the exact special-fibre actions
\begin{equation}\label{p2:eq:two-C3-actions}
 \begin{array}{c|cc}
 &\overline C&E_0\\\hline
 c&e\mapsto\zeta e&(X,Y)\mapsto(\zeta^2X,Y)\\
 \sigma&e\mapsto\zeta e&(X,Y)\mapsto(X,Y).
 \end{array}
\end{equation}
Possible signs from a choice of a square root in the preliminary
quadratic twist disappear on the fourth symmetric power.

The quotient of $\overline C$ by $e\mapsto\zeta e$ is the rational
$v$-line.  Since three is invertible in the $\ell$-adic coefficient
field, finite-map trace gives
\[
 H^1(\overline C,\Q_\ell)^{\sigma}=0.
\]
Thus, by \eqref{p2:eq:two-C3-actions}, the entire main component term
\[
 H^1(\overline C,\Q_\ell)\otimes
       \operatorname{Sym}^4H^1(E_0,\Q_\ell)
\]
has no $\sigma$-invariants.  This remains true for its $A_4$-invariant
part and for every subsequent projector.  The conclusion uses both
the base and coefficient actions, not the coefficient action alone.

\subsubsection{The projective quaternion action on the constant coefficient}
The full $A_4$ action on projective three-torsion remains faithful
after reduction: on a good elliptic model, three-torsion is finite
etale and its four projective lines remain distinct.  Since $K_4$
fixes the main base pointwise, its action on the constant elliptic
coefficient is therefore faithful in $\operatorname{Aut}(E_0)/\{\pm1\}$.
All these are zero-preserving elliptic isomorphisms.  Their projective
ambiguity is exactly the quadratic-twist sign already removed in
$\operatorname{Sym}^4$.

For completeness, the inverse image of this $K_4$ is the quaternion
group $Q_8$.  Its eight automorphisms are
\[
 (X,Y)\longmapsto(X+r,Y+r^2X+t),\qquad
 r\in\mathbf F_4,\quad t^2+t=r^3.
\]
For $r\ne0$ their squares are $[-1]:(X,Y)\mapsto(X,Y+1)$.
The order-three automorphism $(X,Y)\mapsto(\zeta^2X,Y)$ permutes
the three nonidentity elements of the projective $K_4$.  Together
these generate the binary tetrahedral group $Q_8\rtimes C_3$.

Let $M=\operatorname{Sym}^4 H^1(E_0,\Q_\ell)$.  The quaternion
character calculation gives
\begin{equation}\label{p2:eq:Q8-sym4}
 \dim M^{K_4}=\dim M^{Q_8}=2,\qquad
 M^{Q_8\rtimes C_3}=0.
\end{equation}
Indeed the traces on $\operatorname{Sym}^4$ are $5$ on the two
central elements and $1$ on the six elements of order four.  On
each of the sixteen remaining binary tetrahedral elements the
trace is $-1$.  Averaging gives $(10+6)/8=2$ and
$(10+6-16)/24=0$.  More precisely, if $c_0$ is either generator
of the quotient $C_3$, then
\[
 \operatorname{Tr}(c_0\mid M^{Q_8})
 =\frac18\sum_{q\in Q_8}\operatorname{Tr}(c_0q\mid M)=-1.
\]
Together with the absence of fixed vectors this shows that its
characteristic polynomial on $M^{K_4}$ is $X^2+X+1$, so its
eigenvalues are $\zeta,\zeta^2$.
This calculation is valid for every $\ell\ne2$, including $\ell=3$,
because coefficients have characteristic zero.

\subsubsection{Arithmetic action on the two node orbits}
The three attaching points over $z=0$ are $v=0$, $e^3=1$;
the three over $z=\infty$ are the three values of $e/v$ at infinity.
In each case $c$ and $\sigma$ permute these three points in the same
cycle.  The stabilizer in $A_4$ is $K_4$.  Use the good node scale
$k=v^2$ at zero and $k=1/w$ at infinity, where $w=1/v$.
These scales contain neither $e$ nor $\pi$.  Thus $c$ acts trivially
on the displayed elliptic node coordinates, whereas $\sigma$ acts
by $(X,Y)\mapsto(\zeta X,Y)$.

Identifying the $A_4$-invariant direct sum over either three-point
orbit with $M^{K_4}$ at one point, the arithmetic action is therefore
the coefficient action of $c^{-1}\sigma$.  It is the order-three
rotation $(X,Y)\mapsto(\zeta X,Y)$.  By \eqref{p2:eq:Q8-sym4}, neither
of these two two-dimensional node terms has arithmetic invariants.
Their duals and Tate twists have the same property, since the
$\ell$-adic Tate character is unramified at two.

At each of the remaining three main attaching points $e=0$, the
whole geometric $A_4$ is the stabilizer.  Its projective action is
the full tetrahedral action above, so its node coefficient invariant
space is zero by \eqref{p2:eq:Q8-sym4}.  The main degree-zero constant
coefficient invariant space is zero for the same reason.

Thus every main and node term surviving geometric descent has no
arithmetic eigenvalue one. The tail calculation and the exact
cohomological passage below complete the invariant calculation.
\subsection{The descended ordinary tails}
The degree-four tail quotient can be read from the exact rational map.
At zero, $v=\pi u$ gives
\[
\frac{z}{\pi^4}=\frac{u(u^3+1)}{4u^3+1}
\quad\longmapsto\quad Z=u^4+u.
\]
This is the translation quotient by $K_4=(\F_4,+)$. Since
$\lambda=z^3=16Z^3$, the reduced Legendre $j$-invariant is $1/Z^6$.
Infinity is its inversion counterpart.

At a tail over a cube root of unity, use $e=\pi u$ and $h=2u^3-1$.
The exact identity
\[
\lambda-1=\frac{(h-1)^3(h+1)}{(2h+1)^3}
\]
gives
\[
T=\frac{\lambda-1}{16}
=\frac{u^3(u^3-1)^3}{(4u^3-1)^3}
\quad\longmapsto\quad (u^4+u)^3.
\]
This is the degree-twelve quotient by the affine group
$A_4=(\F_4,+)\rtimes\F_4^\times$. Its reduced $j$-invariant is $1/T^2$.
The ordinary characteristic-two family
\[
\mathcal K_t:\quad y^2+xy=x^3+tx
\]
has $j=1/t^2$. Hence the two descended fourth-symmetric-power systems
are respectively its cubic pullback $t=Z^3$ and its ordinary version
$t=T$. Ordinary elliptic curves with the same $j$ differ by a quadratic
twist; this ambiguity disappears on $\Sym^4$.
Over the ordinary generic point, the zero-preserving automorphism group
is $\{\pm1\}$. Thus the identification is compatible with the descended
fourth-symmetric-power action, independently of the choice of elliptic
isomorphism.

These identifications include middle extensions at the missing points.
Indeed finite pushforward on the compact quotient curves is exact, and
finite-group invariants over $\Q_\ell$ are exact. The invariant pushforward
of a middle-extension coefficient system is the middle extension of its
generic descent, as may also be checked on the local inertia-invariant
stalks. This applies to the wildly ramified special translation quotients
as well: it uses characteristic-zero coefficients, not tameness of those
finite maps. Thus the component parabolic cohomology after deck descent
is exactly that of the displayed ordinary or cubic Kloosterman system.
The cited Kloosterman calculation gives zero for both of these parabolic
$\Sym^4$ groups, accounting for every descended tail middle term.
\subsection{Vanishing of the descended tail terms}

\begin{lemma}[The characteristic-two Kloosterman tails]
\label{p2:lem:p2-kl-tails}
Let $k=\overline{\mathbf F}_2$, let $\ell\ne2$, and let
 $f:E\to\mathbf G_{m,k}$ be the elliptic scheme
\[
 E_t:\quad y^2+xy=x^3+tx.
\]
Put $\mathcal F=\operatorname{Sym}^4R^1f_*\mathbf Q_\ell$, and
 let $[3]:\mathbf G_m\to\mathbf G_m$ send $Z$ to $Z^3$.
Then
\[
 H^1_!(\mathbf G_m,\mathcal F)=0,
 \qquad H^1_!(\mathbf G_m,[3]^*\mathcal F)=0.
\]
For the first sheaf the Swan conductor at infinity and the
 dimensions of the inertia-invariant spaces at $0,\infty$ are
 $(1,1,0)$; for the second they are $(3,1,2)$.
\end{lemma}

\begin{proof}
On the open set $x\ne0$, put $u=y/x$.  The elliptic equation becomes
 $u^2+u=x+t/x$, the Artin--Schreier cover defining the rank-two
 Kloosterman sheaf $\mathrm{Kl}_2$ in characteristic two.
The involution $u\mapsto u+1$ is elliptic inversion.  The two
 omitted points $x=0$ and the point at infinity are fixed by this
 involution, so their boundary cohomology is invariant.  Taking
 the odd part of the compact-support exact sequence identifies
 $R^1f_*\mathbf Q_\ell$ with $\mathrm{Kl}_2$.  This also fixes
 the trace convention: its Frobenius trace is minus the usual
 Kloosterman sum.  The elliptic realization is recorded in
 Haessig~\cite{Haessig}, proof of Theorem~2.2, citing Schoof--van der Vlugt,
 Proposition~3.1(4).

We use the characteristic-two local calculation of Yun~\cite{Yun},
 \S4.3.1 and Lemmas~4.3.3--4.3.4.  The image of inertia at
 infinity on $\mathrm{Kl}_2$ is the binary tetrahedral group
 $2.A_4$, with lower ramification filtration
\[
 D_0=2.A_4,\quad D_1=Q_8,\quad D_2=D_3=\{\pm1\},
 \quad D_4=1.
\]
In the fourth symmetric power this gives
 $\operatorname{Swan}_\infty(\mathcal F)=1$ and
 $\mathcal F^{I_\infty}=0$.  At zero, the local monodromy is a
 single unipotent block, so $\dim\mathcal F^{I_0}=1$.
The global geometric monodromy is $\mathrm{SL}_2$; therefore
 $H^0(\mathbf G_m,\mathcal F)=H^2_c(\mathbf G_m,\mathcal F)=0$.
The Grothendieck--Ogg--Shafarevich formula and the boundary
 exact sequence consequently give
\[
 \dim H^1_!(\mathbf G_m,\mathcal F)
 =\operatorname{Swan}_\infty(\mathcal F)
  -\dim\mathcal F^{I_0}-\dim\mathcal F^{I_\infty}=1-1-0=0.
\]

The cubic pullback kills the tame quotient $(2.A_4)/Q_8$ of
 the local image at infinity, leaving image $Q_8$.
On $\operatorname{Sym}^4$, its invariant plane is generated by
 $x^4+y^4$ and $x^2y^2$.  Tame pullback multiplies the Swan
 conductor by three.  Thus the corresponding local data are
 $\operatorname{Swan}_\infty=3$, invariant dimension $1$ at zero,
 and invariant dimension $2$ at infinity.  A finite-index
 restriction still has geometric Zariski closure $\mathrm{SL}_2$,
 so the global $H^0$ and $H^2_c$ again vanish.  The same formula
 now reads $3-1-2=0$.
\end{proof}

\begin{remark}
The lemma applies to the ordinary tail $t=(\lambda-1)/16$ and
 to the cubic tails $\lambda=16Z^3$ supplied by the explicit
 Legendre coordinate changes, after quadratic twists are removed
 by $\operatorname{Sym}^4$.  The statement concerns parabolic
 cohomology of these characteristic-two families.  Identifying
 their contributions to the original dyadic inertia invariants
 still requires the component and node maps of the semistable
 model; the vanishing alone is not a local-factor calculation.
\end{remark}
\subsection{Equivariant nearby cycles and parabolic descent}

We apply the preceding calculations to the proper marked semistable
model $\mathcal B/\mathcal O_L$ of the actual splitting cover
$\eta^3=4(v^3+1)$. Its main component is $C:e^3=v^3+1$, where
$e=\eta/\pi^2$. Write $G=A_4$ for the geometric deck group,
$c:e\mapsto\zeta e$ for the indicated geometric order-three operation,
and $\sigma(\pi)=\zeta\pi$ for the arithmetic descent operation.
The two operations on the coefficient are those computed in the
arithmetic-action calculation above; they must not be identified.

The marked model has a genus-one main component and nine rational tails,
each with four marked cusp sections. It is a stable marked curve.
Consequently the deck and arithmetic descent operations extend to it
by uniqueness of the stable marked model. Let $\mathcal D$ be the union
of the horizontal cusp sections, and write
\[
 j:\mathcal B\setminus\mathcal D\hookrightarrow\mathcal B,
 \qquad
 \mathcal F=\operatorname{Sym}^4R^1f_*\mathbf Q_\ell,
 \qquad \mathcal J=j_*\mathcal F.
\]
Here $f$ is the good elliptic family on the indicated open, supplied by
the integral charts. We assert equivariance for $\mathcal F$, not for
an arbitrary choice of the preliminary quadratic twist of the elliptic
family. On the generic open, $\mathcal F$ is canonically the pullback
of the original fourth symmetric-power coefficient. Its deck and
arithmetic identifications extend to the lisse sheaf on
$\mathcal B\setminus\mathcal D$: restriction of lisse sheaves to a
dense open of a normal scheme is fully faithful. They then extend to
$\mathcal J$ by functoriality of $j_*$. In particular the following
sequences carry the actual $G$ and arithmetic actions.

\subsubsection{The local nearby-cycle calculation}

Let $\mathcal B_a$ denote the normalized special-fibre components, and
let $\mathcal J_a$ be their coefficient middle extensions across their
horizontal cusps. For each of the nine vertical nodes $x$, put
$H_x=\mathcal F_x$ and set
\[
 \mathcal N=\bigoplus_{x\text{ a vertical node}}H_x.
\]
The good elliptic charts at the nodes are essential here. The
semistable-curve calculation gives
\begin{equation}\label{p2:eq:complete-p2-nearby}
 R^0\Psi(\mathcal J_{\bar\eta})=\mathcal J_{\bar s},\qquad
 R^1\Psi(\mathcal J_{\bar\eta})
     =\bigoplus_x(H_x(-1))_x,\qquad
 R^q\Psi(\mathcal J_{\bar\eta})=0\quad(q>1).
\end{equation}
At a vertical node the coefficient is lisse on the full ring
$uv=\pi$. Thus the assertion is the constant-sheaf node calculation
tensored with $H_x$; see \cite[Corollary~8.4.4]{Illusie} for that
semistable specialization of the nearby-cycle formula.

There are no extra vertical vanishing cycles at the horizontal cusps.
The base is smooth there, the marked section is a relative smooth
divisor, and the elliptic family is multiplicative with local parameter
$q=t^6$ times a unit. After strict henselization the branches split and
the unit has every prime-to-$2$ root. The coefficient therefore has its
usual tame unipotent monodromy in $t$, and its horizontal invariant
sheaf $\mathcal F^{I_t}$ is lisse over the trait.
More precisely, \cite[Expose~XIII, Lemmas~2.1.10--2.1.11]{SGA7}
gives base change for $Rj_*\mathcal F$ and
$R\Phi(j_!\mathcal F)=0$ for a lisse sheaf tamely ramified along a
relative normal-crossings divisor. Apply also smooth local acyclicity
to the invariant sheaf on the section. The exact sequence
\[
 0\longrightarrow j_!\mathcal F\longrightarrow j_*\mathcal F
 \longrightarrow i_*\mathcal F^{I_t}\longrightarrow0
\]
then gives $R\Phi(j_*\mathcal F)=0$ at that section. The statements
apply at finite prime-to-$2$ coefficient levels and then to
$\mathbf Q_\ell$ coefficients. Away from the cusps and vertical nodes,
ordinary smooth local acyclicity applies. These observations prove
\eqref{p2:eq:complete-p2-nearby} and identify the restrictions of
$\mathcal J_{\bar s}$ with the component middle extensions.

Proper base change and the five-term nearby-cycle sequence yield
\begin{equation}\label{p2:eq:complete-p2-five-term}
 0\longrightarrow H^1(\mathcal B_{\bar s},\mathcal J_{\bar s})
 \longrightarrow H^1(\mathcal B_{\bar\eta},\mathcal J_{\bar\eta})
 \longrightarrow\mathcal N(-1)
 \longrightarrow H^2(\mathcal B_{\bar s},\mathcal J_{\bar s}).
\end{equation}
For the proper comparison, see
\cite[Expose~XIII, \S2.1.8, especially (2.1.8.9)]{SGA7}.
We only use that the quotient of the second term by the first embeds
in $\mathcal N(-1)$; surjectivity onto the node module is unnecessary.

For normalization $\nu:\coprod_a\mathcal B_a\to\mathcal B_{\bar s}$,
lissity at the nodes gives
\[
 0\longrightarrow\mathcal J_{\bar s}
 \longrightarrow\nu_*\bigoplus_a\mathcal J_a
 \longrightarrow\bigoplus_x(H_x)_x\longrightarrow0.
\]
The last map is the difference of the two restrictions. Orient every
edge from the unique genus-one main component toward its tail. Both
$G$ and $\sigma$ preserve this orientation, so there is no additional
sign in their action on the node module. The cohomology sequence gives
\begin{equation}\label{p2:eq:complete-p2-normalization}
 \bigoplus_a H^0(\mathcal B_a,\mathcal J_a)
 \longrightarrow\mathcal N
 \longrightarrow H^1(\mathcal B_{\bar s},\mathcal J_{\bar s})
 \longrightarrow\bigoplus_a H^1(\mathcal B_a,\mathcal J_a)
 \longrightarrow0.
\end{equation}
Both sequences are equivariant. Taking $G$-invariants is exact over
$\mathbf Q_\ell$, also when $\ell=3$.

\subsubsection{Finite descent, including the wild special-tail quotients}

We recall why taking the indicated stabilizer invariants on a special
tail computes its descended middle cohomology even when the quotient
is wildly ramified at the attaching point. Let
$q:Y\to X=Y/H$ be any of the finite special-tail quotient maps.
On the dense open where the coefficient descends, let $F_X$ be the
descended sheaf and let $j_Y,j_X$ be the respective open immersions.
Finite pushforward and exactness of rational $H$-invariants give
\[
 \bigl(q_*j_{Y*}q^*F_X\bigr)^H=j_{X*}F_X.
\]
This can be checked directly on sections: extension by $j_*$ and
finite pushforward commute, and invariant sections on the finite
cover are precisely sections downstairs. Hence
\[
 H^1(Y,j_{Y*}q^*F_X)^H=H^1(X,j_{X*}F_X).
\]
No tameness of $q$ at the omitted point is required. In the present
case the maps on the open tails are the displayed separable quotients
$Z=u^4+u$ and $T=(u^4+u)^3$. Their ramification at infinity is thus
handled by the middle extension, not discarded. The same argument
justifies generic finite Galois descent of parabolic cohomology.

The preceding component and group calculations now supply exactly the
following inputs:
\begin{enumerate}
\item The main-component cohomology has no $\sigma$-eigenvalue $1$:
$\sigma$ acts on its Fermat-base $H^1$ with the two nontrivial cubic
characters and acts trivially on the constant elliptic coefficient.
\item All descended tail middle-cohomology groups vanish, by their
identification with the ordinary and cubic-pullback Kloosterman
families and the stated local-monodromy calculation.
\item $\mathcal N^G$ consists of the two two-dimensional node spaces
over $0$ and $\infty$. Each has arithmetic eigenvalues
$\zeta,\zeta^2$; the three other node spaces vanish after $G$-invariants.
The Tate twist $(-1)$ is unramified at $2$.
\end{enumerate}
Consequently neither $\mathcal N^G$ nor
$(\bigoplus_a H^1(\mathcal B_a,\mathcal J_a))^G$ has eigenvalue $1$.
Equation~\eqref{p2:eq:complete-p2-normalization} gives the same conclusion
for $H^1(\mathcal B_{\bar s},\mathcal J_{\bar s})^G$.
Equation~\eqref{p2:eq:complete-p2-five-term} then gives it for
$H^1(\mathcal B_{\bar\eta},\mathcal J_{\bar\eta})^G$.
Absence of an eigenvalue passes to subspaces, quotients and extensions,
so no degeneration assertion or vanishing of graph monodromy is used.

Finally, $H^1$ of the coefficient middle extension on a proper smooth
curve is canonically its parabolic cohomology. Generic finite Galois
descent therefore identifies the last $G$-invariant group with the full
original cubic-Legendre parabolic representation. An inertia-invariant
vector would be fixed by the chosen arithmetic inertia lift $\sigma$,
which has no eigenvalue $1$ here. Both specified rank-four summands
thus have zero inertia invariants. Combining this with Corollary~\ref{p2:cor:true-cover-swan} gives, for every $\ell\ne2$,
\[
 d_{2,\ell}^+=d_{2,\ell}^-=0,\qquad
 P^+_{2,\ell}(T)=P^-_{2,\ell}(T)=1,\qquad
 a_2(W^+_{4,\ell})=a_2(W^-_{4,\ell})=4.
\]


This proves Theorem~\ref{p2:thm:dyadic-local-data}. In particular the
arithmetic constant in the splitting cover is retained throughout;
replacing $\eta^3=4(v^3+1)$ by a geometrically isomorphic equation
without tracking that constant would not determine this inertia action.
